\section{Exact Weil heights within the rational isogeny class}
\label{sec:uniform-weil-height}

The complex absolute value in
Theorem~\ref{thm:entire-class-bounds} is a single local quantity.
We now determine the absolute Weil height of the same four
$j$-invariants, including every cancellation in their rational
coordinates. The minimum over the entire rational isogeny class is
$3\log(c^2-16c+16)$. In particular its leading term remains
$6\log c$, although the logarithm of the minimum complex absolute
value has leading term $\log c$.

\subsection{Reduced rational coordinates}

Throughout this section put
\begin{equation}\label{eq:wh-parameters}
 n\in\Z_{\ge1},\qquad c=1792n+2,\qquad u=c/2=896n+1.
\end{equation}
Thus $c\ge1794>32$, and $u,c-1$ are positive odd integers.
Let $E_c,E_0,E_a,E_b$ be the four actual curves of
Theorem~\ref{thm:entire-four-models}, and write
\begin{equation}\label{eq:wh-polynomials}
 \begin{aligned}
 P&=c^2-c+1,& Q&=c^2-16c+16,\\
 R&=c^2+14c+1,& S&=16c^2-16c+1.
 \end{aligned}
\end{equation}
For a rational number $z=N/D$ in lowest terms, with $D>0$, we use
the absolute logarithmic and multiplicative Weil heights
\begin{equation}\label{eq:wh-rational-height}
 \height(z)=\log\max\{|N|,D\},\qquad
 \mathrm H(z)=\max\{|N|,D\}.
\end{equation}
Over $\Q$ the relative and absolute normalizations agree.

\begin{proposition}\label{prop:wh-reduced-fractions}
The following are the signed reduced rational coordinates of the
actual $j$-invariants. Every denominator is positive.
\[
\begin{array}{c|r|l}
 \text{curve}&\text{numerator }N&\text{denominator }D\\ \hline
 E_c&64P^3&u^2(c-1)^2\\
 E_0&-Q^3&(c-1)u^4\\
 E_a&8R^3&u(c-1)^4\\
 E_b&8S^3&u(c-1)
\end{array}
\]
\end{proposition}

\begin{proof}
The previously computed invariants give
\[
 \begin{aligned}
 j(E_c)&=\frac{256P^3}{c^2(c-1)^2},&
 j(E_0)&=-\frac{16Q^3}{(c-1)c^4},\\
 j(E_a)&=\frac{16R^3}{c(c-1)^4},&
 j(E_b)&=\frac{16S^3}{c(c-1)}.
 \end{aligned}
\]
Substitution of $c=2u$ gives the displayed coordinates.
The sign of $j(E_0)$ follows from its negative displayed
discriminant. This computation uses $j=c_4^3/\Delta$ and does not
require the displayed models to be minimal at $2$.

For coprimality, the exact residue identities give the table
\[
\begin{array}{c|cc}
 &\bmod u&\bmod(c-1)\\ \hline
 P&1&1\\
 Q&16&1\\
 R&1&16\\
 S&1&1
\end{array}
\]
For example,
$Q=16+u(4u-32)=1+(c-1)(c-15)$ and
$R=1+u(4u+28)=16+(c-1)(c+15)$.
Because $u,c-1$ are odd, $16$ is coprime to either modulus.
Consequently each polynomial is coprime to both factors of every
denominator. Taking powers and multiplying numerator factors
$64$ or $8$ preserves this coprimality.

For completeness the exact $2$-parts are also visible directly:
$P,R,S$ are odd, whereas
\[
 Q=4(u^2-8u+4),\qquad v_2(Q)=2.
\]
The $2$-adic valuations of the four absolute numerators are
respectively $6,6,3,3$, and every denominator is odd.
Thus no odd prime or power of $2$ can cancel in any of the
four displayed fractions.
\end{proof}

\subsection{The two exact minima}

Let $\mathcal I_c$ denote the rational isogeny class of $E_c$,
modulo rational isomorphism, and define
\[
 \height_{\min}(c)=\min_{F\in\mathcal I_c}\height(j(F)),\qquad
 J_{\min,\infty}(c)=\min_{F\in\mathcal I_c}|j(F)|_\infty.
\]
Theorem~\ref{thm:entire-four-models} supplies the entire-class
identification used in these minima; the following arithmetic
calculations do not provide a replacement proof of that
classification.

\begin{lemma}\label{lem:wh-complex-minimum}
The unique minimizing rational isomorphism class for the complex
absolute value is represented by $E_0$, and
\begin{equation}\label{eq:wh-complex-minimum}
 J_{\min,\infty}(c)
   =\frac{Q^3}{(c-1)u^4}
   =\frac{16Q^3}{(c-1)c^4}>2c>1.
\end{equation}
In particular all four absolute numerators in
Proposition~\ref{prop:wh-reduced-fractions} exceed their denominators.
\end{lemma}

\begin{proof}
For $c\ge32$ one has
\begin{equation}\label{eq:wh-polynomial-inequalities}
 \begin{aligned}
 Q-\frac{c^2}{2}&=\frac{c(c-32)}2+16>0,& Q&<c^2,\\
 P-Q&=15(c-1)>0,\\
 R-Q&=15(2c-1)>0,\\
 S-Q&=15(c^2-1)>0.
 \end{aligned}
\end{equation}
Thus all four polynomials are positive.
Since $Q>c^2/2$ and $c-1<c$, the absolute value of $j(E_0)$
is greater than $2c$. The exact ratios of the other three
absolute values to this one are
\begin{equation}\label{eq:wh-three-ratios}
 \begin{aligned}
 \frac{|j(E_c)|}{|j(E_0)|}
   &=\frac{16c^2}{c-1}\left(\frac P Q\right)^3>1,\\
 \frac{|j(E_a)|}{|j(E_0)|}
   &=\left(\frac{cR}{(c-1)Q}\right)^3>1,\\
 \frac{|j(E_b)|}{|j(E_0)|}
   &=\left(\frac{cS}{Q}\right)^3>1.
 \end{aligned}
\end{equation}
These inequalities follow from
\eqref{eq:wh-polynomial-inequalities} and $c>1$.
The entire-class theorem now gives the assertion, including
uniqueness. Strictness separates $E_0$ from all other displayed
models, regardless of any possible coincidence among those others.
\end{proof}

\begin{theorem}\label{thm:wh-exact-minimum}
The absolute logarithmic Weil heights of the actual four
$j$-invariants are
\begin{equation}\label{eq:wh-height-table}
 \begin{aligned}
 \height(j(E_c))&=\log64+3\log P,\\
 \height(j(E_0))&=3\log Q,\\
 \height(j(E_a))&=\log8+3\log R,\\
 \height(j(E_b))&=\log8+3\log S.
 \end{aligned}
\end{equation}
The unique minimizer in the entire rational isogeny class is
represented by $E_0$, with
\begin{equation}\label{eq:wh-exact-minimum}
 \height_{\min}(c)=3\log(c^2-16c+16),\qquad
 \min_{F\in\mathcal I_c}\mathrm H(j(F))=Q^3.
\end{equation}
Moreover,
\begin{equation}\label{eq:wh-height-bounds}
 \begin{aligned}
 6\log c-3\log2&<\height_{\min}(c)<6\log c,\\
 c^6/8&<\min_{F\in\mathcal I_c}\mathrm H(j(F))<c^6.
 \end{aligned}
\end{equation}
\end{theorem}

\begin{proof}
The fractions in Proposition~\ref{prop:wh-reduced-fractions} are
reduced, and Lemma~\ref{lem:wh-complex-minimum} shows that their
absolute numerators exceed their denominators. Therefore
\eqref{eq:wh-rational-height} gives exactly
\eqref{eq:wh-height-table}.
The inequalities $P,R,S>Q>0$, together with $64,8>1$, show that
each height other than that of $E_0$ is strictly greater than
$3\log Q$. The entire-class theorem extends this minimum and
uniqueness to every rationally isogenous curve.
Finally $c^2/2<Q<c^2$ gives \eqref{eq:wh-height-bounds}, first
by cubing and then by taking logarithms.
\end{proof}

\subsection{Finite-place contribution and the bounded gain}

The minimizers above are the same, but the two minimum values
measure different quantities. If $N/D$ is reduced, each prime
dividing $D$ contributes exactly $v_q(D)\log q$ to the finite-place
sum in the logarithmic Weil height; primes not dividing $D$
contribute zero. The entire finite-place contribution is therefore
$\log D$. For $E_0$, whose complex absolute value exceeds one,
this proves the exact identity
\begin{equation}\label{eq:wh-finite-place-gap}
 \begin{aligned}
 \height_{\min}(c)-\log J_{\min,\infty}(c)
    &=\log((c-1)u^4)\\
    &=5\log c-\log16+\log(1-1/c).
 \end{aligned}
\end{equation}
Equivalently, the exact formulas yield
\begin{equation}\label{eq:wh-limits}
 \begin{aligned}
 \height_{\min}(c)-6\log c
    &=3\log(1-16/c+16/c^2)\longrightarrow0,\\
 \frac{J_{\min,\infty}(c)}c
    &=16\,\frac{(1-16/c+16/c^2)^3}{1-1/c}
       \longrightarrow16.
 \end{aligned}
\end{equation}
All limits here and below are along the unbounded sequence
\eqref{eq:wh-parameters}. Continuity of $\log$ at positive arguments
proves these limits. In particular
\[
 \frac{\height_{\min}(c)}{\log c}\longrightarrow6,\qquad
 \frac{\log J_{\min,\infty}(c)}{\log c}\longrightarrow1.
\]
The denominator accounts for the missing term of size $5\log c$;
it cannot be discarded when using the absolute Weil height.

The gain achieved by choosing the best representative in the
entire class is also exact:
\begin{equation}\label{eq:wh-bounded-gain}
 \height(j(E_c))-\height_{\min}(c)
   =\log64+3\log(P/Q)\in(6\log2,9\log2).
\end{equation}
Indeed $P>Q$ by \eqref{eq:wh-polynomial-inequalities}, while
$P<c^2$ and $Q>c^2/2$ imply $P/Q<2$. These strict inequalities
prove the stated interval, using $\log64=6\log2$.
Since $P/Q\to1$, the gain tends to $6\log2$.
Thus optimization over the entire class saves only a bounded
additive amount in this family.

\begin{corollary}\label{cor:wh-no-leading-exponent-saving}
For every $\delta>0$ and $C\in\R$, there exists an $n\ge1$ such
that every elliptic curve $F/\Q$ rationally isogenous to $E_{c_n}$
satisfies
\[
 \height(j(F))>(6-\delta)\log c_n+C.
\]
Equivalently, for every $\theta<6$ and $C_0>0$, some $n\ge1$
satisfies
\[
 \mathrm H(j(F))>C_0c_n^\theta
 \qquad\text{for every }F\in\mathcal I_{c_n}.
\]
\end{corollary}

\begin{proof}
Choose $n$ with $\delta\log c_n>C+3\log2$ and apply the lower
bound in \eqref{eq:wh-height-bounds}. Such $n$ exist because
$c_n\to\infty$. For the multiplicative formulation, choose
$c_n^{6-\theta}>8C_0$ and use the corresponding lower bound.
\end{proof}

This corollary excludes the specific proposed step of lowering the
leading sixth-power $j$-height size by selecting a different rationally
isogenous representative. It is not a counterexample to a
conductor--height estimate, modified Szpiro, or abc. No upper bound
for the radicals of these triples is asserted, and their unresolved
height--radical comparison is unaffected by this calculation.

\subsection{Formalization boundary}

The Lean height companion reuses the four actual Weierstrass curves
from the preceding section.
It proves the signed rational equalities, integer coprimality, and
equality with the canonical \texttt{Rat.num} and \texttt{Rat.den}
coordinates. It then uses the actual multiplicative and logarithmic
height functions in mathlib and proves compatibility with the existing
normalized-height interface over $\Q$. This gives the exact height
tables, attained minima over the actual model enumeration,
uniqueness, and the logarithmic bounds
\eqref{eq:wh-height-bounds}. No artificially defined height is
substituted for the library functions.

The enumeration is not asserted in Lean to exhaust a rational
isogeny class. That passage uses the separately proved
Theorem~\ref{thm:entire-four-models}, including its stated external
isogeny and Frobenius inputs. The strict comparison of the complex
minimum, the explicit finite-place identity, the asymptotic limits,
the bounded-gain formula, and
Corollary~\ref{cor:wh-no-leading-exponent-saving} are paper proofs
in this section; they are not added to the companion's formal
claims. The module adds no axioms or conjectural hypotheses.
